Soit $ABC$ un triangle et $J$ le centre de son cercle exinscrit opposé au sommet $A$. Ce cercle est tangent au
côté $[BC]$ en $M$ et aux droites $(AB)$ et $(AC)$, respectivement, en $K$ et $L$. Les droites $(LM)$ et
$(BJ)$ se coupent en $F$ et les droites $(KM)$ et $(CJ)$ se coupent en $G$. Soit $S$ le point d'intersection
des droites $(AF)$ et $(BC)$ et soit $T$ le point d'intersection des droites $(AG)$ et $(BC)$.

Montrer que $M$ est le milieu du segment $[ST]$.

(Le cercle \emph{exinscrit} du triangle $ABC$ opposé au sommet $A$ est le cercle tangent au segment $[BC]$, à
la demi-droite $[AB]$ au-delà de $B$ et à la demi-droite $[AC]$ au-delà de $C$)
